%%% FIELD joint-spectral-limit-proof
Let \(Z=(y,M,B,s_1)\), where the matrix blocks are vectorized in a fixed order, and let \(\widehat Z=(\widehat y,\widehat M,\widehat B,\widehat s_1)\). By Assumption~\ref{ass:fixed-panel}, \(Z\) lies in a Euclidean space of fixed dimension. Lemma~\ref{lem:joint-cell-expansion} gives the stacked expansion
\[
 \widehat Z-Z=K^{-1}\sum_{i=1}^K\Phi_i+r_K,
 \qquad \sqrt K\,\|r_K\|_2\to_P0.
\]
The definition of each centered cell influence vector gives \(E\Phi_i=0\). Assumptions~\ref{ass:cell-iid} and~\ref{ass:cross-cell-independence} imply that the stacked vectors \(\Phi_1,\ldots,\Phi_K\) are independent and identically distributed: within a fixed \(i\), their cell blocks are independent, and across distinct \(i\) all constituent observations are independent. Because the number and dimensions of the blocks are fixed, Assumption~\ref{ass:influence-moment} implies
\[
 E\|\Phi_i\|_2^{2+\nu}<\infty.
\]
Consequently
\[
 E\left\|K^{-1/2}\sum_{i=1}^K\Phi_i\right\|_2^2
 =E\|\Phi_i\|_2^2<\infty,
\]
where the cross terms vanish by independence and centering. Markov's inequality and the displayed expansion therefore yield
\[
 \|\widehat Z-Z\|_2=O_P(K^{-1/2}),
 \qquad
 \|\widehat M-M\|_{\mathrm{op}}=O_P(K^{-1/2}).
\]

We next verify that no separate singular-gap assumption is needed. Assumption~\ref{ass:exact-rank} and \(r\ge1\) imply
\[
 \sigma_r(M)>0,
 \qquad
 \sigma_{r+1}(M)=0,
\]
where the second equality uses the convention \(\sigma_{r+1}(M)=0\) if \(r=\min(N_0,T_0)\). Since \(M\) is the fixed population donor-pre matrix and \(K\) is the sole asymptotic index, \(\kappa:=\sigma_r(M)/3\) is a fixed positive number. The variational characterization of singular values gives, for every conformable \(A\),
\[
 |\sigma_j(A)-\sigma_j(M)|\le \|A-M\|_{\mathrm{op}}.
\]
Thus, whenever \(\|A-M\|_{\mathrm{op}}<\kappa\),
\[
 \sigma_r(A)>2\kappa,
 \qquad
 \sigma_{r+1}(A)<\kappa,
\]
so the leading rank-\(r\) singular subspace is separated on this neighborhood. The preceding stochastic bound shows that \(\widehat M\) belongs to this neighborhood with probability tending to one.

Assumption~\ref{ass:exact-rank} also gives \(M_r=M\) and \(P=M^+\), with
\[
 \|P\|_{\mathrm{op}}=\sigma_r(M)^{-1}<\infty.
\]
On the separated neighborhood just derived, Lemma~\ref{lem:full-spectral-differential} establishes Fr\'echet differentiability of \(\mathcal F(y,M,B)=yM_r^+B\), with the derivative recorded in Definition~\ref{def:full-differential}. Hence there is a remainder \(\rho_K\) such that
\[
 \frac{\|\rho_K\|_2}{\|\widehat Z-Z\|_2}\to_P0
\]
and, by Definition~\ref{def:full-spectral-map},
\[
\begin{aligned}
 \widehat\theta-\theta
 &=(\widehat y-y)PB
   +y\mathcal L_M(\widehat M-M)B
   +yP(\widehat B-B)+\rho_K\\
 &=(\widehat y-y)PB
   +y\{-P(\widehat M-M)P
      +PP^\top(\widehat M-M)^\top(I-MP)\\
 &\hspace{4.9cm}
      +(I-PM)(\widehat M-M)^\top P^\top P\}B
   +yP(\widehat B-B)+\rho_K.
\end{aligned}
\]
Since \(\|\widehat Z-Z\|_2=O_P(K^{-1/2})\), the defining property of \(\rho_K\) gives \(\rho_K=o_P(K^{-1/2})\). Substituting the target-pre, donor-pre, and donor-post blocks of Lemma~\ref{lem:joint-cell-expansion} into the first display above gives
\[
\begin{aligned}
 \widehat\theta-\theta
 &=K^{-1}\sum_{i=1}^K
 \{\Phi_{y,i}PB+y\mathcal L_M(\Phi_{M,i})B+yP\Phi_{B,i}\}
 +o_P(K^{-1/2})\\
 &=K^{-1}\sum_{i=1}^K\Phi_{\theta,i}+o_P(K^{-1/2}),
\end{aligned}
\]
where the final equality is Definition~\ref{def:transported-influence}. The factual-post block of Lemma~\ref{lem:joint-cell-expansion} is
\[
 \widehat s_1-s_1
 =K^{-1}\sum_{i=1}^K\Phi_{1,i}+o_P(K^{-1/2}).
\]
Stacking the last two equations and again invoking Definition~\ref{def:transported-influence} proves
\[
 \sqrt K
 \begin{pmatrix}
  \widehat\theta-\theta\\
  \widehat s_1-s_1
 \end{pmatrix}
 =K^{-1/2}\sum_{i=1}^K\Phi_{\eta,i}+o_P(1).
\]
This calculation explicitly retains perturbations from \(\widehat y\), \(\widehat M\), and \(\widehat B\), including both projector-correction terms in \(\mathcal L_M\).

The map from \(\Phi_i\) to \(\Phi_{\eta,i}\) in Definition~\ref{def:transported-influence} is linear and fixed. Its coefficients are finite because \(y\) and \(B\) are finite-dimensional survival arrays, \(P\) is bounded, and the displayed formula for \(\mathcal L_M\) is a bounded linear map. Therefore
\[
 E\Phi_{\eta,i}=0,
 \qquad
 E\|\Phi_{\eta,i}\|_2^{2+\nu}<\infty,
 \qquad
 E(\Phi_{\eta,i}\Phi_{\eta,i}^\top)=\Omega_\eta.
\]
For completeness, the Gaussian limit can be obtained without nonsingularity. Fix \(a\in\mathbb R^{2J}\), let
\[
 W_i=a^\top\Phi_{\eta,i},
 \qquad
 \omega_a^2=a^\top\Omega_\eta a,
 \qquad
 \delta=\min(\nu,1)>0.
\]
The preceding moment bound implies \(E|W_i|^{2+\delta}<\infty\), while \(EW_i=0\) and \(EW_i^2=\omega_a^2\). Taylor expansion on \(|x|\le1\), together with the boundedness of the left side on \(|x|>1\) after division by \(|x|^{2+\delta}\), supplies a finite constant \(C_\delta\) satisfying
\[
 \left|e^{\mathrm i x}-1-\mathrm i x+\frac{x^2}{2}\right|
 \le C_\delta |x|^{2+\delta}
 \quad\text{for every }x\in\mathbb R.
\]
It follows that
\[
 E e^{\mathrm iW_i/\sqrt K}
 =1-\frac{\omega_a^2}{2K}+o(K^{-1}).
\]
Independence then gives
\[
 E\exp\!\left\{\mathrm i a^\top K^{-1/2}
                   \sum_{i=1}^K\Phi_{\eta,i}\right\}
 =\left(1-\frac{\omega_a^2}{2K}+o(K^{-1})\right)^K
 \longrightarrow e^{-\omega_a^2/2}.
\]
The right side is the characteristic function of the scalar projection of \(\mathcal N_{2J}(0,\Omega_\eta)\). Since this holds for every \(a\), the characteristic-function criterion for weak convergence yields
\[
 K^{-1/2}\sum_{i=1}^K\Phi_{\eta,i}
 \rightsquigarrow \mathcal N_{2J}(0,\Omega_\eta).
\]
If \(a^\top\Omega_\eta a=0\), the corresponding limit projection is the point mass at zero; hence the argument permits singular \(\Omega_\eta\). Combining this limit with the asymptotic linear representation proves the asserted joint convergence.

It remains to establish the two plug-in conclusions. For a matrix \(A\) in the separated neighborhood, let \(\Pi_R(A)\) be the orthogonal projector onto its leading \(r\) right singular subspace and put \(A_r=A\Pi_R(A)\). The gap derived above makes \(\Pi_R(A)\) unique. To see continuity at \(M\), take any sequence \(A_m\to M\). The set of rank-\(r\) orthogonal projectors is compact. Every convergent subsequence of \(\Pi_R(A_m)\) maximizes
\[
 Q\longmapsto \operatorname{tr}(QM^\top M)
\]
over that set, by continuity and the variational characterization. Because \(\sigma_r(M)>\sigma_{r+1}(M)\), this maximizer is uniquely \(\Pi_R(M)\). Thus every convergent subsequence has the same limit, and
\[
 \Pi_R(A)\to\Pi_R(M)
 \quad\text{as }A\to M.
\]
Moreover,
\[
 A_r^+=\{A_r^\top A_r+I-\Pi_R(A)\}^{-1}A_r^\top.
\]
Indeed, on the range of \(\Pi_R(A)\) the matrix in braces has eigenvalues \(\sigma_1(A)^2,\ldots,\sigma_r(A)^2\), while on its orthogonal complement it is the identity. These eigenvalues are bounded away from zero in the neighborhood. Continuity of multiplication and inversion therefore gives
\[
 (\widehat M)_r^+\to_PP.
\]
Substitution in the explicit polynomial formula in Definition~\ref{def:full-differential}, together with \(\widehat M\to_PM\), yields convergence of the corresponding differential maps,
\[
 \|\widehat{\mathcal L}_{\widehat M}-\mathcal L_M\|_{\mathrm{op}}\to_P0.
\]

Let \(\widehat\Phi_i\) denote the stack of the fold-specific plug-in cell influence vectors. The final assertion of Lemma~\ref{lem:joint-cell-expansion} states foldwise \(L_2(P_c)\)-consistency. Conditional on each training sample, the validation observations are independent of the fitted nuisance functions. Hence conditional Markov inequalities, followed by a union bound over the fixed numbers of cells, folds, and coordinates, give
\[
 K^{-1}\sum_{i=1}^K\|\widehat\Phi_i-\Phi_i\|_2^2=o_P(1).
\]
Let \(\mathsf A\) be the fixed linear map that sends \(\Phi_i\) to \(\Phi_{\eta,i}\), and let \(\widehat{\mathsf A}\) be its plug-in analogue. Consistency of \(\widehat y,\widehat M,\widehat B\), the last two spectral convergences, and Definition~\ref{def:transported-influence} imply
\[
 \|\widehat{\mathsf A}-\mathsf A\|_{\mathrm{op}}=o_P(1),
 \qquad
 \|\widehat{\mathsf A}\|_{\mathrm{op}}=O_P(1).
\]
The weak law and \(E\|\Phi_i\|_2^2<\infty\) give \(K^{-1}\sum_i\|\Phi_i\|_2^2=O_P(1)\). Therefore
\[
\begin{aligned}
 K^{-1}\sum_{i=1}^K
 \|\widehat\Phi_{\eta,i}-\Phi_{\eta,i}\|_2^2
 &\le
 2\|\widehat{\mathsf A}\|_{\mathrm{op}}^2
 K^{-1}\sum_{i=1}^K\|\widehat\Phi_i-\Phi_i\|_2^2\\
 &\quad+
 2\|\widehat{\mathsf A}-\mathsf A\|_{\mathrm{op}}^2
 K^{-1}\sum_{i=1}^K\|\Phi_i\|_2^2
 =o_P(1).
\end{aligned}
\]

Write \(D_i=\widehat\Phi_{\eta,i}-\Phi_{\eta,i}\). Expanding the difference of outer products and applying Cauchy--Schwarz gives
\[
\begin{aligned}
 &\left\|\widehat\Omega_\eta-
 K^{-1}\sum_{i=1}^K\Phi_{\eta,i}\Phi_{\eta,i}^\top\right\|_{\mathrm{op}}\\
 &\quad\le
 \left(K^{-1}\sum_{i=1}^K\|D_i\|_2^2\right)^{1/2}
 \left\{
  \left(K^{-1}\sum_{i=1}^K\|\widehat\Phi_{\eta,i}\|_2^2\right)^{1/2}
  +
  \left(K^{-1}\sum_{i=1}^K\|\Phi_{\eta,i}\|_2^2\right)^{1/2}
 \right\}
 =o_P(1).
\end{aligned}
\]
Every coordinate product of \(\Phi_{\eta,i}\) is integrable by Cauchy--Schwarz. The scalar weak law applied entrywise, with fixed dimension \(2J\), consequently gives
\[
 K^{-1}\sum_{i=1}^K\Phi_{\eta,i}\Phi_{\eta,i}^\top
 \to_P\Omega_\eta.
\]
The last two displays prove \(\widehat\Omega_\eta\to_P\Omega_\eta\), without positive definiteness.

Finally, the probability limit furnished by Lemma~\ref{lem:joint-cell-expansion} is the observed-law cell functional
\[
 S_c(t)=E_{P_c}\{Q_c(t\mid X_{ci})\},
\]
whose identification under covariate-dependent censoring uses Assumptions~\ref{ass:conditional-censoring} and~\ref{ass:censor-positivity}; it is not a marginal Kaplan--Meier pseudo-target. The symbol \(\theta=yM_r^+B\) names an observed-law spectral functional. Assumption~\ref{ass:post-transport}, rather than that definition or any row-space condition on \(y\), equates this functional to the target unit's untreated potential post-period survival row. The row \(s_1\) is the factual treated post-period survival row. This proves the claimed causal interpretation and completes the proof.
